\documentclass[a4paper,12pt]{article}
\usepackage[utf8]{inputenc}
\usepackage{graphicx}
\usepackage[serbian]{babel}
\usepackage{hyperref}
\usepackage{tocloft}

\begin{document}

\begin{titlepage}
    \centering
\begin{figure}[htbp]
    \centering
    \includegraphics[width=0.4\textwidth]{./logo.png}
\end{figure}
    { Univerzitet u Beogradu \\ Matematički fakultet\par}

    \vfill

    {\Large \textbf{Projekat iz Računarske inteligencija}\par}

    \vspace{1cm}

    {\Large \textbf{Primena genetskog algoritma u proceni nedostajućih vrednosti u randomizovanom kompletnom blok dizajnu 
}\par}

    \vfill

   


\begin{tabbing}
\hspace{10cm} \= \hspace{10cm} \= \kill
\textbf{Mentor:} \>  \textbf{Studenti:} \\
Prof. dr Vladimir Filipović \> Dragana Katić 91/2021 \\
Asistent Stefan Kapunac \> Đurđa Milošević 84/2021 \\
\> Smer: Informatika
\end{tabbing}

    \vfill

    \textbf{Datum:} avgust 2025

\end{titlepage}

\newpage

\tableofcontents

\section{Uvod}


\section{Teorijska osnova}
\subsection{Randomizovani kompletni blok dizajn (RCBD)}
\subsection{Genetski algoritam (GA)}


\section{Opis rešenja}
\subsection{Opis problema}

\section{Poređenje sa drugim metodama}
\subsection{Poređenje sa Eq. 3, prosečna vrednost}
\subsection{Grafički prikaz konvergencije}
\subsection{Diskusija rezultata}

\section{Demonstracija metode}
\subsection{}
\subsection{}
\subsection{}


\section{Zaključak}
\begin{itemize}
\item Postignuti rezultati
\item Prednosti i ograničenja GA pristupa
\item Moguće unapređenje i budući rad
\end{itemize}

\section{Literatura}

\end{document}
